\section{Conclusion}

Crop-ranking reversal is a conditional allocation phenomenon generated by
stochastic margins, joint
downside risk and operational feasibility. A genuine agronomic score remains
separate from profit. A full loss-CVaR-constrained programme then determines
whether the high-ranked crop receives more or less land.

The restricted two-crop theorem provides an exact frontier rather than an
informal risk-adjusted ranking. When the higher-ranked crop also has the
higher mean margin, reversal occurs precisely when the risk-feasible interval
ends below the equal-share boundary. Unique dependence thresholds require
named-family monotonicity and crossing; otherwise reversal regions can be
nonmonotone or disconnected. The complete KKT system and optimal-face bounds
prevent local multipliers or one solver vertex from being overinterpreted.

In the public-data-calibrated Kansas experiment, winter wheat ranks first in
historical relative-yield potential but receives the least land in the
full model. This is a selected complete rank reversal, not a strong
exclusion reversal, and it is primarily margin-induced; CVaR moderates it.
The principal strong-reversal count is structurally zero because positive crop
minimums prevent exclusion. Complete all-zero and top-crop-zero sensitivity
phases allow crop exit but retain a substantive strong-reversal null.
Controlled experiments separately exhibit risk- and operationally induced
Soybean--Corn crossings. A full-investment matched-Kendall Gaussian
Selected Gaussian mean--variance policy chosen by a fixed 15\%
variance-reduction target
strictly reduces variance relative to its expected-profit benchmark yet
violates the Student-\(t\) loss-CVaR ceiling. The failure persists over a
non-singleton \(\gamma\) interval and is parameter-dependent across one-factor
sensitivity cases. Information--
flexibility interaction in the shared-CVaR model is reported from signed
adjacent-grid cross-differences; the strict theorem remains restricted to a
separable or nonbinding-risk submodel.

Heuristic benchmarks require a declared feasibility projection. Euclidean and
\(L^1\) projection agree for the already feasible proportional and equal-share
rules but give different score-winner allocations and reversal
classifications. That benchmark result is projection-sensitive; the optimized
model results are not altered by the projection choice.

The \StateYears{}-state-year panel shows that score and acreage leaders often
differ, but its leakage-free lagged estimate includes zero. Aggregate
prevalence therefore cannot identify private risk, constraints or welfare.
Future empirical work needs farm-level pre-plant beliefs, local prices and
costs, crop histories, contracts, credit, capacity and signal receipt.

The final message is operational. Rankings can inform planting, but only a
transparent decision model can prescribe acreage. That model should report
its joint tail law, risk frontier, active constraints, multiple-optimum
classification and the action pathway through which information can matter.
